\subsection{Millikans Öltröpfchenexperiment: Bestimmung der Elementarladung}
\label{subsec:millikan-oeltroepfchenexperiment}

Nach Thomsons Experimenten war klar:
Kathodenstrahlen bestehen aus negativ geladenen Teilchen.
Außerdem konnte Joseph John Thomson\index{Thomson, Joseph John}
das Verhältnis von Ladung zu Masse bestimmen,

\[
\frac{e}{m_e}.
\]

Damit war ein entscheidender Schritt getan.
Aber eine wichtige Frage blieb offen:

\begin{quote}
	Wie groß ist die elektrische Ladung eines einzelnen Elektrons?
\end{quote}

Erst wenn diese Ladung bekannt war,
konnte man aus dem Verhältnis \(\frac{e}{m_e}\)
auch die Masse des Elektrons\index{Elektronenmasse} berechnen.

\begin{HistoryBox}[Robert Andrews Millikan]
	\label{box:millikan}
	Robert Andrews Millikan\index{Millikan, Robert Andrews} führte ab 1909
	seine berühmten Öltröpfchenexperimente\index{Öltröpfchenexperiment} durch.
	Dabei gelang es ihm,
	die elektrische Elementarladung\index{Elementarladung} mit hoher Genauigkeit
	zu bestimmen. Für diese Arbeiten erhielt er 1923 den Nobelpreis für Physik.
\end{HistoryBox}

\newpage
\noindent
\subsubsection{Die Grundidee des Experiments}
\label{subsubsec:grundidee-millikan}

Millikan verwendete winzige Öltröpfchen,
die zwischen zwei horizontalen Metallplatten schwebten.
Die Platten bildeten einen Kondensator\index{Kondensator}.
Zwischen ihnen konnte ein elektrisches Feld\index{Elektrisches Feld}
erzeugt werden.

Die Öltröpfchen wurden durch eine feine Düse in den Raum zwischen den
Platten eingebracht. Einige dieser Tröpfchen erhielten dabei elektrische
Ladung, etwa durch Reibung oder durch Ionisation der Luft.

\FloatBarrier
\begin{figure}[H]
	\centering
	\includegraphics[width=\textwidth]{bilder/millikan_oeltroepfchen.png}
	\caption{Schematischer Versuchsaufbau von Millikans Öltröpfchenexperiment: Geladene Öltröpfchen befinden sich zwischen zwei Kondensatorplatten und können durch ein elektrisches Feld in der Schwebe gehalten werden.}
	\label{fig:millikan-oeltroepfchenexperiment}
\end{figure}

Auf jedes geladene Öltröpfchen wirkten mehrere Kräfte:

\begin{itemize}
	\item die Gewichtskraft\index{Gewichtskraft} nach unten,
	\item die elektrische Kraft\index{Elektrische Kraft} im elektrischen Feld,
	\item der Luftwiderstand\index{Luftwiderstand} bei der Bewegung des Tröpfchens.
\end{itemize}

Die zentrale Idee war einfach, aber genial:
Wenn die elektrische Kraft die Gewichtskraft genau ausgleicht,
kann ein geladenes Öltröpfchen in der Schwebe gehalten werden.

\FloatBarrier
\begin{PhysicsBox}[Kräftegleichgewicht beim Öltröpfchen]
	\label{box:kraeftegleichgewicht-oeltroepfchen}
	Schwebt ein geladenes Öltröpfchen zwischen den Kondensatorplatten,
	dann gilt näherungsweise:
	
	\[
	qE = m g.
	\]
	
	Die elektrische Kraft \(qE\) nach oben gleicht die Gewichtskraft
	\(mg\) nach unten aus.
\end{PhysicsBox}

Hierbei ist \(q\) die Ladung des Öltröpfchens,
\(E\) die elektrische Feldstärke,
\(m\) die Masse des Tröpfchens
und \(g\) die Erdbeschleunigung.

Das Kräftegleichgewicht\index{Kräftegleichgewicht} erlaubte es,
aus bekannten Größen auf die Ladung des Tröpfchens zu schließen.
Damit wurde elektrische Ladung nicht mehr nur qualitativ beschrieben,
sondern quantitativ messbar.

\subsubsection{Warum Öltröpfchen?}
\label{subsubsec:warum-oeltroepfchen}

Millikan verwendete Öl,
weil Öltröpfchen sehr klein und zugleich relativ stabil sind.
Wassertröpfchen würden zu schnell verdunsten und dadurch ihre Masse verändern.
Öltröpfchen bleiben dagegen länger erhalten
und konnten unter dem Mikroskop beobachtet werden.

Das war entscheidend:
Nur wenn Größe und Bewegung eines Tröpfchens ausreichend genau bestimmt
werden konnten,
ließ sich seine Ladung berechnen.

\begin{DidacticBox}[Der Trick des Experiments]
	\label{box:trick-millikan}
	Millikan maß nicht direkt die Ladung eines einzelnen Elektrons.
	Er maß die Ladung vieler einzelner Öltröpfchen.
	Dann erkannte er:
	Alle gemessenen Ladungen waren Vielfache einer kleinsten Ladung.
	Diese kleinste Ladung ist die Elementarladung.
\end{DidacticBox}

\subsubsection{Quantisierung der Ladung}
\label{subsubsec:quantisierung-der-ladung}

Millikans entscheidende Erkenntnis bestand nicht allein darin,
die Ladung eines einzelnen Öltröpfchens zu bestimmen.
Wichtiger war der Vergleich vieler Messungen.
Historisch war Millikans Experiment nicht nur eine Messung eines Zahlenwertes.
Um 1900 stand auch die Frage im Raum,
ob elektrische Ladung kontinuierlich veränderbar ist
oder ob sie aus kleinsten Einheiten besteht.
Millikan griff damit eine zentrale Idee der Elektronentheorie auf:
Wenn das Elektron ein reales Teilchen ist,
dann sollte seine Ladung eine feste, kleinste Einheit darstellen.

Sein Öltröpfchenexperiment zeigte genau dieses Muster.
Die gemessenen Ladungen der Tröpfchen waren nicht beliebig,
sondern traten als ganzzahlige Vielfache eines kleinsten Wertes auf.
Damit wurde die Quantisierung der elektrischen Ladung experimentell greifbar.
Für verschiedene geladene Tröpfchen erhielt man unterschiedliche Ladungswerte.
Diese Werte waren jedoch nicht beliebig verteilt.
Sie ließen sich vielmehr als ganzzahlige Vielfache eines kleinsten gemeinsamen
Wertes schreiben:

\[
q = n e,
\qquad n = 1,2,3,\dots
\]

Dabei bezeichnet \(e\) die Elementarladung.

Anschaulich bedeutet dies:
Ein Öltröpfchen konnte die Ladung \(e\), \(2e\), \(3e\), \(4e\) usw. tragen,
aber keine beliebigen Zwischenwerte.
Die elektrische Ladung verhielt sich also nicht wie eine kontinuierlich
veränderbare Größe,
sondern wie eine Größe mit kleinsten Bausteinen.

Millikan fand diese kleinste Einheit nicht durch eine einzelne Messung,
sondern durch das Muster vieler Messwerte.
Wenn die gemessenen Ladungen immer wieder als Vielfache desselben Wertes
auftreten,
dann liegt nahe,
dass dieser Wert die kleinste Ladungseinheit ist.

\begin{MathBox}[Ganzzahlige Vielfache der Elementarladung]
	\label{box:ganzzahlige-vielfache-elementarladung}
	\small
	Die elektrische Ladung\index{Quantisierung der Ladung} eines Körpers
	tritt in ganzzahligen Vielfachen der Elementarladung auf:
	\[
	q = n e,
	\qquad n \in \mathbb{Z}.
	\]
	Für makroskopische Körper ist \(n\) meist sehr groß,
	sodass Ladung scheinbar kontinuierlich wirkt.
	Auf mikroskopischer Ebene zeigt sich jedoch ihre diskrete Struktur.
\end{MathBox}

Der entscheidende Befund war daher nicht nur der Zahlenwert

\[
e \approx 1{,}602 \cdot 10^{-19}\,\mathrm{C},
\]

sondern der Nachweis,
dass elektrische Ladung gequantelt ist.
Damit wurde das Elektron nicht nur als Teilchen mit Masse,
sondern auch als Träger einer fundamentalen elektrischen Ladung verstanden.

\subsubsection{Vom Verhältnis \texorpdfstring{\(e/m_e\)}{e/me} zur Elektronenmasse}
\label{subsubsec:verhaeltnis-em-zur-elektronenmasse}

Millikans Ergebnis war besonders wichtig,
weil es mit Thomsons Ergebnis kombiniert werden konnte.

Thomson hatte das Verhältnis

\[
\frac{e}{m_e}
\]

bestimmt.
Millikan bestimmte nun \(e\).
Damit ließ sich die Masse des Elektrons berechnen:

\[
m_e = \frac{e}{e/m_e}.
\]

Mit anderen Worten:
Aus dem bekannten Ladungs-Masse-Verhältnis und der bekannten
Elementarladung folgt die Masse des Elektrons.

\begin{PhysicsBox}[Warum Millikans Experiment so wichtig war]
	\label{box:bedeutung-millikan}
	Thomson zeigte,
	dass es ein negativ geladenes Teilchen mit einem bestimmten Verhältnis
	von Ladung zu Masse gibt.
	Millikan bestimmte die Ladung dieses Teilchens.
	Zusammen ergab sich daraus die Masse des Elektrons.
\end{PhysicsBox}

Damit wurde das Elektron von einer experimentellen Vermutung
zu einem messbaren physikalischen Objekt mit konkreten Eigenschaften:

\[
q_e = -e,
\qquad
m_e \approx 9{,}109 \cdot 10^{-31}\,\mathrm{kg}.
\]

\subsubsection{Bedeutung für die Physik}
\label{subsubsec:bedeutung-millikan-fuer-physik}

Millikans Öltröpfchenexperiment lieferte einen direkten experimentellen
Hinweis darauf,
dass elektrische Ladung gequantelt ist.
Elektrische Ladung tritt nicht beliebig fein abgestuft auf,
sondern in Vielfachen einer kleinsten Einheit.

Diese Erkenntnis war grundlegend für die moderne Atomphysik.
Sie zeigte,
dass Elektronen nicht nur Träger negativer Ladung sind,
sondern dass ihre Ladung eine fundamentale Naturkonstante darstellt.

\begin{NoteBox}[Einordnung im roten Faden]
	\label{box:einordnung-millikan}
	Millikans Experiment schließt eine wichtige Lücke:
	Nach Thomson kannte man das Verhältnis von Ladung zu Masse.
	Durch Millikan kannte man die Elementarladung.
	Damit wurden Ladung und Masse des Elektrons quantitativ bestimmbar.
\end{NoteBox}